\begin{abstract}
Energy estimates for the three-dimensional incompressible Navier--Stokes equations bound the
kinetic energy, but they do not stop velocity gradients from concentrating in ever smaller regions.
This is why the classical regularity criteria, such as Beale--Kato--Majda \citep{bkm_1984}, ask for
control of a peak quantity. This paper builds a multiscale \emph{decision scaffold}, inspired by
the Six Birds framework, around a different question: can activity cascade to arbitrarily fine
scales in finite time, a \emph{Zeno cascade}? The scaffold keeps three kinds of statement apart.

\emph{Proved in Lean.} If every refinement step costs a fixed quantum of work, and the work
available per unit time at depth $j$ is bounded by a capacity whose reciprocals have divergent sum,
then no Zeno cascade occurs. Route-stable reductions have at most polynomial capacity growth.
In a finite model, a uniform capacity certificate over inputs that include localized tests
forces a pointwise bound on the gain.

\emph{Proved on paper, or stated as open hypotheses.} We define the subgrid bridge with sharp
Fourier shells and show that on a strong solution it returns the actual closure force. Energy
alone bounds the power that the bridge supplies to shell $j$ by $C\,2^{5j/2}\|u_0\|_2^3$. That
bound is too weak to certify divergence. An explicit three-mode field shows that no bound of the
form $\|\nabla u\|_\infty\|u_j\|_2^2$ holds for this power.

\emph{Measured.} An anti-localization certificate on initial data passes band-limited random fields
($\max_j\kappa_j\approx 2.1$) and rejects a localized shell wavepacket ($\kappa\approx 404$).

We do not prove the Clay regularity statement, and we do not show that the certificate persists
under the flow. The Navier--Stokes inequalities that a proof along these lines would still need
are listed explicitly.
\end{abstract}

\noindent\textbf{Keywords:} Navier--Stokes; multiscale scaffold; No-Zeno engine; route capacity;
anti-localization; peak-channel regularity; Lean; Six Birds framework
